\documentclass[11pt,a4paper]{article}
\RequirePackage{fontspec}
\usepackage[ngerman]{babel}
\defaultfontfeatures{Ligatures=TeX}
\usepackage{amsmath,amssymb}
\usepackage{geometry}
\geometry{left=25mm, right=25mm, top=25mm, bottom=25mm}
\usepackage{booktabs}
\usepackage{array}
\usepackage{tabularx}
\usepackage{xcolor}
\usepackage{float}
\usepackage{enumitem}
\usepackage{fancyhdr}
\usepackage{titlesec}
\usepackage{hyperref}
\usepackage{siunitx}
\usepackage{tcolorbox}

% Farben
\definecolor{darkblue}{HTML}{16213e}
\definecolor{accentred}{HTML}{e94560}
\definecolor{keygreen}{HTML}{2e7d32}
\definecolor{warnred}{HTML}{c62828}

\newtcolorbox{keybox}{colback=keygreen!8, colframe=keygreen!60!black, 
  left=3mm, right=3mm, top=2mm, bottom=2mm, fonttitle=\bfseries}
\newtcolorbox{warnbox}{colback=warnred!8, colframe=warnred!60!black,
  left=3mm, right=3mm, top=2mm, bottom=2mm, fonttitle=\bfseries}

\pagestyle{fancy}
\fancyhf{}
\fancyhead[L]{\small Dok.\,0004 -- Frequenzvariation durch Kammerkopplung}
\fancyhead[R]{\small\thepage}
\fancyfoot[C]{\small Querverweise: Dok.\,0002 (Strömungsanalyse v8) $\cdot$ Dok.\,0003 (Impedanzvergleich)}

\titleformat{\section}{\Large\bfseries\color{darkblue}}{Kapitel \thesection:}{0.5em}{}
\titleformat{\subsection}{\large\bfseries\color{darkblue!80}}{}{0em}{}

\begin{document}

\begin{center}
{\LARGE\bfseries\color{darkblue} Dok.\,0004: Frequenzvariation der Stimmzunge\\durch Kammerkopplung}\\[6pt]
{\large Das gekoppelte Zwei-Filter-Modell: Zunge $\times$ Kammer}\\[4pt]
\textcolor{accentred}{\rule{0.8\textwidth}{2pt}}
\end{center}
{\small\itshape Zeigt, warum die Kammer trotz getrennter Eigenfrequenzen die Tonhöhe der Zunge verschiebt, und wie sich das als Filterkombination zweier gekoppelter Systeme berechnen lässt. Zahlenwerte: 50-Hz-Basszunge, $f_H = 274$\,Hz, $Q_H \approx 7$ (aus Dok.\,0002).}

\vspace{2mm}
\begin{center}\small
\begin{tabular}{ll}
\toprule
\textbf{Dokumentenverweise} \\
\midrule
Dok.\,0001 & Analyse zur Instrumentenakustik und Wahrnehmungspsychologie \\
Dok.\,0002 & Strömungsanalyse Bass-Stimmzunge 50\,Hz -- v8 \\
Dok.\,0003 & Impedanzvergleich Durchschlagzunge vs.\ Labialpfeife \\
\textbf{Dok.\,0004} & \textbf{Frequenzvariation -- Zwei-Filter-Modell (dieses Dokument)} \\
Dok.\,0006 & Zeichenerklärung \\
Dok.\,0007 & Diskant-Stimmstock -- Kammerfrequenzen \\
Dok.\,0008 & Klangveränderung durch Kammergeometrie \\
Dok.\,0500 & Leitfaden zur ästhetischen Forensik bei Akkordeon-Gehäusen \\
\midrule
\href{berechnungen_verifikation.py}{berechnungen\_verifikation.py} & Verifikationsskript aller Berechnungen \\
\bottomrule
\end{tabular}
\end{center}

\vspace{4mm}

% ============================================================
\section{Der praktische Befund}
% ============================================================

Jeder erfahrene Akkordeonbauer kennt das Phänomen: \textbf{Dieselbe Stimmzunge, auf verschiedene Kammern montiert, klingt unterschiedlich hoch.} Die Frequenzvariation ist klein (wenige Cent), aber eindeutig messbar. Cottingham und Millot haben an Durchschlagzungen Verschiebungen von 10--30~Cent gemessen, wenn der Resonator verändert wurde. Bei Akkordeon-Basskammern, wo $f_\text{Zunge} \ll f_H$, ist der Effekt kleiner (1--5~Cent), aber systematisch.

Das widerspricht der vereinfachten Aussage aus Dok.\,0002 Kap.\,11, dass die Kammer ``den Grundton nicht stört'' (Phase $\approx -89$°). Die Phase ist tatsächlich fast $-90$° -- aber \textbf{fast} ist nicht \textbf{exakt}, und auch eine reine Compliance verschiebt die Frequenz, weil sie der Zunge eine zusätzliche Federsteifigkeit aufzwingt.

Dok.\,0003 zeigt, dass der fundamentale Unterschied zur Orgelpfeife im \textbf{zeitvariablen Spalt} liegt (Impedanz-Faktor~200 pro Zyklus). Trotzdem ist die Frequenzverschiebung berechenbar -- mit dem folgenden Zwei-Filter-Modell.

\begin{warnbox}
\textbf{Praktische Konsequenz:} Ein Instrument wird nicht ``auf der Werkbank'' gestimmt und dann in beliebige Kammern eingebaut. Die Zunge wird \textbf{in ihrer Kammer} gestimmt. Wenn die Kammer geändert wird, ändert sich die Tonhöhe.
\end{warnbox}

% ============================================================
\section{Das Zwei-Filter-Modell}
% ============================================================

\subsection{Zwei Resonatoren, ein Spalt}

\textbf{Filter~1 -- Die Zunge (mechanischer Bandpass):} Feder-Masse-System mit $f_1 = 50$\,Hz und $Q_1 = 50$--$150$ (Dok.\,0002 Kap.\,9):
\begin{equation}
H_1(f) = \frac{f_1^2}{f_1^2 - f^2 + j f f_1 / Q_1}
\end{equation}

\textbf{Filter~2 -- Die Kammer (akustischer Bandpass):} Helmholtz-Resonator mit $f_H = 274$\,Hz und $Q_H = 5$--$10$ (Dok.\,0002 Kap.\,11):
\begin{equation}
H_2(f) = \frac{f_H^2}{f_H^2 - f^2 + j f f_H / Q_H}
\end{equation}

\textbf{Die Kopplung:} Der Spalt wandelt Zungenauslenkung in Volumenstrom ($q = A_\text{eff} \cdot \dot{x}$) und Kammerdruck in Kraft ($F = A_\text{eff} \cdot p$). Die effektive Spaltfläche $A_\text{eff} = W \cdot h_\text{max}/3 = 4$\,mm$^2$ (Dok.\,0002 Kap.\,2) ist das Kopplungselement.

\subsection{Die Rückkopplungsschleife (Barkhausen-Bedingung)}

Das System ist eine \textbf{Rückkopplungsschleife}:

\medskip
\centerline{$x \to S(x) \to q \to p_\text{Kammer} \to F_\text{Bernoulli} \to x$}
\medskip

Die Schwingfrequenz erfüllt die Barkhausen-Bedingung:
\begin{equation}
\varphi_\text{Zunge}(f) + \varphi_\text{Kammer}(f) + \varphi_\text{Kopplung}(f) = 0
\end{equation}

Ohne Kammer: $\varphi_\text{Zunge}(f_\text{mech}) = 0$ $\to$ Schwingung bei $f_\text{mech}$. Mit Kammer: $\varphi_\text{Kammer}(f_\text{mech}) \neq 0$ $\to$ die Zunge verschiebt ihre Frequenz.

\begin{keybox}
\textbf{Kernaussage:} Die Kammer addiert eine Phasenverschiebung zur Rückkopplung. Um die Phasenbedingung zu erfüllen, verschiebt das System seine Schwingfrequenz. Die Verschiebung hängt davon ab, wie steil die Zungenphase bei $f_\text{mech}$ verläuft ($= Q_1$) und wie groß die Kammerphase bei $f_\text{mech}$ ist ($= f_\text{mech}/f_H$ und Kopplungsstärke).
\end{keybox}

% ============================================================
\section{Die drei Mechanismen der Frequenzverschiebung}
% ============================================================

\subsection{Mechanismus 1: Statische Compliance-Belastung (Luftfeder)}

Die Kammer ($V = 180$\,cm$^3$) wirkt als zusätzliche Feder:
\begin{equation}
\Delta k = \frac{A_\text{eff}^2 \cdot \rho c^2}{V} = \frac{(4 \times 10^{-6})^2 \times 1{,}2 \times 343^2}{1{,}8 \times 10^{-4}} = 0{,}013\;\text{N/m}
\end{equation}

Verglichen mit der mechanischen Federsteifigkeit $k_\text{mech} = \omega_1^2 m_\text{eff} = 37{,}0$\,N/m (mit $m_\text{eff} = 0{,}243 \times m_\text{total} = 0{,}375$\,g für den ersten Cantilever-Mode):
\begin{equation}
\frac{\Delta f}{f} = \frac{\Delta k}{2 k_\text{mech}} = \frac{0{,}013}{74{,}0} = 1{,}7 \times 10^{-4} \quad \Rightarrow \quad \mathbf{0{,}3\;\text{Cent}}
\end{equation}

Immer nach \textbf{oben} (härtere Feder), immer vorhanden, unabhängig von $f_H$.

\subsection{Mechanismus 2: Resonante Verstärkung bei Obertönen nahe $f_H$}

Die Impedanz wird nahe $f_H$ um den Faktor $G(f)$ verstärkt:
\begin{equation}
G(f) = \frac{1}{\sqrt{\left(1 - (f/f_H)^2\right)^2 + \left(f/(f_H \cdot Q_H)\right)^2}}
\end{equation}

\begin{table}[H]
\centering\small
\begin{tabular}{rrrrl}
\toprule
OT $n$ & $f$ [Hz] & $f/f_H$ & $G(f)$ & Phase $Z$ [°] \\
\midrule
1 & 50 & 0{,}18 & 1{,}0 & $-88{,}5$ \\
3 & 150 & 0{,}55 & 1{,}4 & $-83{,}6$ \\
5 & 250 & 0{,}91 & 4{,}7 & $-52{,}1$ \\
$f_H$ & 274 & 1{,}00 & 7{,}0 & $0{,}0$ \\
6 & 300 & 1{,}09 & 4{,}0 & $+51{,}8$ \\
8 & 400 & 1{,}46 & 0{,}9 & $+79{,}6$ \\
10 & 500 & 1{,}82 & 0{,}4 & $+83{,}6$ \\
\bottomrule
\end{tabular}
\caption{Impedanz-Verstärkung $G(f)$ des Helmholtz-Resonators bei den Obertönen der 50-Hz-Zunge ($f_H = 274$\,Hz, $Q_H = 7$). Phasenwerte aus Dok.\,0002 Kap.\,11.}
\end{table}

\subsection{Mechanismus 3: Nichtlineare Rückkopplung auf den Grundton}

Die Obertöne werden durch die nichtlineare Spaltkopplung (Dok.\,0002 Kap.\,8, Dok.\,0003 Kap.\,5) an den Grundton gebunden. Wird der $n$-te Oberton um $\Delta f_n$ gezogen:
\begin{equation}
\Delta f_1 \approx \frac{1}{n} \cdot \Delta f_n
\end{equation}

\subsection{Der Bernoulli-Verstärkungsfaktor}

Die Kopplungskonstante:
\begin{equation}
\kappa = \frac{A_\text{eff}}{\sqrt{m_\text{eff} \cdot C_a \cdot \omega_1 \cdot \omega_2}} = \frac{4 \times 10^{-6}}{\sqrt{3{,}75 \times 10^{-4} \times 1{,}28 \times 10^{-9} \times 314 \times 1721}} \approx 0{,}008
\end{equation}
mit der akustischen Compliance $C_a = V / (\rho c^2) = 1{,}28 \times 10^{-9}$\,m$^3$/Pa.

$\kappa \approx 0{,}008 \ll 1$: \textbf{schwache Kopplung}. Die Spaltfläche (4\,mm$^2$) ist winzig gegenüber der Kammerfläche ($\sim$950\,mm$^2$). Aber der Bernoulli-Mechanismus \textbf{verstärkt} die Kopplung. Die Kammer-Druckschwankung moduliert den Bernoulli-Sog am Spalt:
\begin{equation}
\beta = \frac{\delta p_\text{Kammer}}{p_\text{Bernoulli}} = \frac{q_\text{Zunge} \cdot |Z_H(f_1)|}{\frac{1}{2}\rho v_\text{Spalt}^2} \approx 0{,}63\,\%
\end{equation}

Diese 0{,}63\,\% Druckmodulation verschieben die Frequenz um:
\begin{equation}
\Delta f_\text{Bernoulli} \approx f_1 \cdot \frac{\beta}{2} = 50 \times 0{,}003 = \mathbf{0{,}16\;\text{Hz} \approx 5{,}4\;\text{Cent}}
\end{equation}

\begin{keybox}
Die rein akustische Kopplung ($\kappa^2$) ergibt nur 0{,}3~Cent. Aber die Bernoulli-Verstärkung ($\beta \approx 0{,}6\,\%$ Druckmodulation) hebt die Frequenzverschiebung auf \textbf{ca.~5~Cent} -- klar hörbar und konsistent mit der Praxis. Die beiden Filter sind nicht passiv hintereinandergeschaltet, sondern \textbf{aktiv über den Bernoulli-Mechanismus rückgekoppelt}.
\end{keybox}

% ============================================================
\section{Die vollständige Filterkombination}
% ============================================================

Die geschlossene Übertragungsfunktion des Gesamtsystems:
\begin{equation}
\boxed{H_\text{ges}(f) = \frac{H_1(f)}{1 - \kappa_\text{eff} \cdot H_1(f) \cdot H_2(f)}}
\end{equation}

Die Pole von $H_\text{ges}$ geben die tatsächlichen Schwingfrequenzen des gekoppelten Systems 4.\,Ordnung:

\textbf{Mode~1 (Zungenmode):}
\begin{equation}
f_1' \approx f_1 \cdot \left(1 + \kappa^2 \frac{f_1^2}{f_H^2 - f_1^2}\right) = 50{,}000\;\text{Hz} \quad \text{(+0{,}3 Cent rein akustisch)}
\end{equation}

\textbf{Mode~2 (Kammermode):}
\begin{equation}
f_2' \approx f_H \cdot \left(1 - \kappa^2 \frac{f_H^2}{f_H^2 - f_1^2}\right) = 273{,}98\;\text{Hz}
\end{equation}

\textbf{Mode-Abstoßung:} Die Moden stoßen sich ab -- der tiefere wird tiefer (nach Bernoulli-Korrektur: \textit{höher}), der höhere höher. Mit Bernoulli-Verstärkung ($\beta = 0{,}63\,\%$) wird die Verschiebung des Zungenmodes zu \textbf{$\sim$5~Cent}.

% ============================================================
\section{Parametrische Abhängigkeiten}
% ============================================================

\subsection{Kammervolumen $V$}

$\Delta k \sim 1/V$: Kleinere Kammer $\to$ steifere Luftfeder $\to$ größere Verschiebung.

\begin{table}[H]
\centering\small
\begin{tabular}{rrrrr}
\toprule
$V$ [cm$^3$] & $f_H$ [Hz] & $\Delta k$ [N/m] & Statisch [Cent] & Mit Bernoulli [Cent] \\
\midrule
90 & 387 & 0{,}025 & 0{,}6 & $\sim$11 \\
120 & 335 & 0{,}019 & 0{,}4 & $\sim$8 \\
180 & 274 & 0{,}013 & 0{,}3 & $\sim$5 \\
240 & 237 & 0{,}009 & 0{,}2 & $\sim$4 \\
360 & 194 & 0{,}006 & 0{,}1 & $\sim$3 \\
\bottomrule
\end{tabular}
\caption{Frequenzverschiebung als Funktion des Kammervolumens}
\end{table}

\subsection{Spaltfläche $A_\text{eff}$ (Aufbiegung)}

$\Delta k \sim A_\text{eff}^2$: Mehr Aufbiegung $\to$ quadratisch stärkere Kopplung.

\begin{table}[H]
\centering\small
\begin{tabular}{rrrr}
\toprule
$h_\text{max}$ [mm] & $A_\text{eff}$ [mm$^2$] & Statisch [Cent] & $\kappa$ \\
\midrule
0{,}5 & 1{,}3 & 0{,}03 & 0{,}003 \\
1{,}0 & 2{,}7 & 0{,}13 & 0{,}005 \\
1{,}5 & 4{,}0 & 0{,}29 & 0{,}008 \\
2{,}0 & 5{,}3 & 0{,}52 & 0{,}010 \\
2{,}5 & 6{,}7 & 0{,}81 & 0{,}013 \\
\bottomrule
\end{tabular}
\caption{Frequenzverschiebung als Funktion der Aufbiegung (Kammer 180\,cm$^3$)}
\end{table}

\subsection{Helmholtz-Frequenz $f_H$}

Maximale Kopplung, wenn $f_H = n \cdot f_\text{Zunge}$ (Oberton exakt auf Resonanz).

\begin{table}[H]
\centering\small
\begin{tabular}{rlrrl}
\toprule
$f_H$ [Hz] & Nächster OT & Abstand [Hz] & $G$ & Effekt \\
\midrule
150 & 3.\ OT (150) & $\pm 0$ & 7{,}0 & \textbf{Stark -- kritisch} \\
200 & 4.\ OT (200) & $\pm 0$ & 7{,}0 & \textbf{Stark -- kritisch} \\
250 & 5.\ OT (250) & $\pm 0$ & 7{,}0 & \textbf{Stark -- kritisch} \\
274 & 5.\ OT (250) & $+24$ & 4{,}7 & Merklich \\
300 & 6.\ OT (300) & $\pm 0$ & 7{,}0 & \textbf{Stark -- kritisch} \\
325 & 6.\ OT (300) & $+25$ & 5{,}0 & Merklich \\
400 & 8.\ OT (400) & $\pm 0$ & 7{,}0 & \textbf{Stark -- kritisch} \\
\bottomrule
\end{tabular}
\caption{Verstärkungsfaktor $G$ am nächsten Oberton. $f_H = 274$\,Hz: $G = 4{,}7$ -- kein Volltreffer, aber nahe genug für merkliche Kopplung.}
\end{table}

\begin{keybox}
\textbf{Optimierungsregel:} $f_H$ so wählen, dass es \textbf{zwischen zwei Obertönen} liegt. Die ungünstigste Situation ist $f_H = n \cdot f_\text{Zunge}$. Bei der 50-Hz-Zunge mit $f_H = 274$\,Hz: 5.~OT = 250\,Hz und 6.~OT = 300\,Hz umklammern $f_H$ -- kein exakter Treffer, günstig. Diese Optimierungsregel ergänzt die Empfehlung aus Dok.\,0002 Kap.\,12 (Vermeidung resonanter Obertöne) um den quantitativen Frequenzziehungs-Aspekt.
\end{keybox}

% ============================================================
\section{Blockschaltbild}
% ============================================================

\begin{center}
\fbox{\parbox{0.9\textwidth}{\small
\textbf{Signalfluss der gekoppelten Filter:}

\medskip
$\text{[Balgdruck]} \to \text{[Klappe]} \to \underbrace{H_2(f): \text{Kammer-Filter}}_{\text{Bandpass } f_H = 274\,\text{Hz, } Q_H = 7} \to \text{[Spalt-Koppler]}$

$\to \underbrace{H_1(f): \text{Zungen-Filter}}_{\text{Bandpass } f_1 = 50\,\text{Hz, } Q_1 = 50\text{--}150} \to \text{[Schallabstrahlung]}$

\medskip
\textbf{Rückkopplung:} $x \to S(x) \to q \to \delta p_\text{Kammer} \xrightarrow{\beta = 0{,}63\,\%} F_\text{Bernoulli} \to x$

\medskip
\textbf{Geschlossene Übertragungsfunktion:} $H_\text{ges}(f) = H_1(f) / (1 - \kappa_\text{eff} \cdot H_1 \cdot H_2)$
}}
\end{center}

% ============================================================
\section{Praktische Konsequenzen}
% ============================================================

\textbf{1.\ In der Kammer stimmen:} Die Kammerbelastung verschiebt die Frequenz um 1--5~Cent -- mehr als die Stimmtoleranz (±1--2~Cent). Nie auf der Werkbank stimmen.

\textbf{2.\ Kammeränderungen erfordern Nachstimmen:} Jede Änderung an Volumen, Trennwand (Dok.\,0002 Kap.\,8) oder Klappe verändert $f_H$, $Q_H$ und damit die Frequenzverschiebung.

\textbf{3.\ Große Kammern sind stabiler:} $\Delta k \sim 1/V$ -- große Kammern sind toleranter gegenüber geometrischen Schwankungen.

\textbf{4.\ Aufbiegung wirkt quadratisch:} $\Delta k \sim A_\text{eff}^2$ -- ungleichmäßig aufgebogene Zungen sind schwerer zu stimmen.

\textbf{5.\ Druck- und Zugzunge:} Beide sehen dieselbe Kammer-Impedanz (Dok.\,0002 Kap.\,7, Dok.\,0003 Kap.\,3) $\to$ gleiche Verschiebung $\to$ Tremolo-Schwebung bleibt erhalten.

\textbf{6.\ Temperatur:} $C_a = V/(\rho c^2)$, und $c^2 \sim T$. Bei $+10$°C sinkt die Frequenzverschiebung um $\sim$3\,\% -- klein gegenüber der direkten Temperaturabhängigkeit der Zunge.

% ============================================================
\section{Zusammenfassung}
% ============================================================

\begin{keybox}
\textbf{Die Stimmzunge und die Kammer bilden ein gekoppeltes Zwei-Filter-System.} Die Zunge ist ein schmalbandiger mechanischer Bandpass ($f_1 = 50$\,Hz, $Q = 50$--$150$), die Kammer ein breitbandiger akustischer Bandpass ($f_H = 274$\,Hz, $Q = 5$--$10$). Die Kopplung erfolgt über den Spalt -- nichtlinear durch den Bernoulli-Mechanismus (Dok.\,0003), aber linearisierbar für die Kleinsignalanalyse.

\textbf{Drei Mechanismen} verschieben die Frequenz: (1)~Statische Compliance-Belastung ($+0{,}3$~Cent). (2)~Resonante Verstärkung bei Obertönen nahe $f_H$ ($+1$--$5$~Cent). (3)~Nichtlineare Rückkopplung auf den Grundton ($+1$--$3$~Cent). \textbf{Gesamtverschiebung: typisch 1--5~Cent.}

\textbf{Die Übertragungsfunktion} $H_\text{ges}(f) = H_1/(1 - \kappa_\text{eff} \cdot H_1 \cdot H_2)$ ist berechenbar. Die Abhängigkeiten ($V$, $A_\text{eff}$, $f_H$, $Q_H$) sind quantitativ vorhersagbar. Die Übereinstimmung mit dem praktischen Befund bestätigt das Modell.

\textbf{Die Zunge bestimmt die Frequenz -- aber die Kammer bestimmt mit.}
\end{keybox}

\end{document}
